% Plot-Level Semantic Similarity Detection in Literary Texts
% Computational Linguistics Journal Submission (MIT Press)
%
% NOTE: This version uses standard article class.
% Once you download clv2025.cls from submissions.cljournal.org/index.php/cljournal/StyleGuide,
% replace the \documentclass line with: \documentclass{clv2025}
% and update bibliography style to: \bibliographystyle{compling}

\documentclass[11pt]{article}
\usepackage[utf8]{inputenc}
\usepackage{times}
\usepackage{latexsym}
\usepackage{amsmath}
\usepackage{graphicx}
\usepackage{booktabs}
\usepackage{algorithm}
\usepackage{algorithmic}
\usepackage{subcaption}
\usepackage{multirow}
\usepackage{natbib}
\usepackage[margin=1in]{geometry}
\usepackage{hyperref}

% Title - 15pt bold, centered, 2.5cm from top (per CL journal style)
% Updated title to distinguish from DH2022 presentation and emphasize benchmark contribution
\title{\Large\bfseries Evaluating Sentence Encoders for Plot-Level Semantic Similarity:\\
A Systematic Comparison of Bi-Encoders, Cross-Encoders, and LLMs}

% Authors - SINGLE-BLIND (not anonymous)
\author{
Grant Glass \\
\\
University of North Carolina at Chapel Hill \\
Department of English and Comparative Literature \\
\texttt{grantg@live.unc.edu}
}

\date{}

\begin{document}

\maketitle

% Abstract: 150-250 words (informative) - Current: ~250 words ✓
\begin{abstract}
Recent embedding models (E5, BGE) dominate sentence-level semantic similarity benchmarks, but whether their performance transfers to document-level compositional semantic tasks remains unexplored. We introduce plot-level semantic similarity detection as an evaluation task that tests whether embeddings capture high-level narrative structure independent of lexical and stylistic features. Using Robinson Crusoe adaptations (1719-1903) as a case study, we construct a balanced binary classification dataset of 2,968 full-length texts where positive examples share core narrative patterns but differ in vocabulary, style, setting, and character names. We systematically evaluate 15 models across three architectural paradigms: bi-encoders (USE, SBERT, E5, BGE), cross-encoders (BERT, RoBERTa), and encoder-decoder models (T5, FLAN-T5), plus five recent LLMs (Llama, Mistral, Phi-3, Gemma, Qwen). Our key findings challenge conventional assumptions: (1) Universal Sentence Encoder (99.3\% accuracy) outperforms modern embeddings (E5: 98.4\%, BGE: 98.2\%) despite inferior STS benchmark scores, revealing inverse correlation between benchmark rankings and our task; (2) cross-encoders provide only 0.2\% gain at 48× computational cost, making them infeasible for corpus-scale analysis; (3) zero-shot LLMs achieve 76-91\%, far below supervised baselines. We release our benchmark dataset, models, and code, and discuss implications for embedding model selection beyond standard benchmarks.
\end{abstract}

\noindent\textbf{Keywords:} compositionality, distributional semantics, sentence embeddings, semantic similarity, document-level semantics, benchmark dataset, evaluation methodology

\section{Introduction}

A central question in computational linguistics is whether distributional semantic models capture \textit{compositional} meaning---the principle that the meaning of a complex expression is determined by the meanings of its parts and how they combine \cite{Partee1995,Fodor2025Compositionality}. While sentence embedding models achieve impressive results on semantic textual similarity (STS) benchmarks \cite{Cer2017SemEval,Muennighoff2023MTEB}, these evaluations primarily test lexical paraphrase at the sentence level. An open question remains: \textit{do these models capture compositional semantics at the document level}, where texts share high-level structure but differ completely in surface features?

We address this question by introducing \textbf{plot-level semantic similarity detection} as a challenging evaluation task. Unlike STS benchmarks that reward lexical overlap, our task requires models to recognize abstract narrative structure despite zero lexical similarity---a strong test of compositional semantic understanding. Consider these examples:

\begin{quote}
\small
\textbf{Text A (1719):} ``I was born in York. After shipwreck, I found myself on a desolate island where I survived alone for years...''

\textbf{Text B (1854):} ``Captain Grant's vessel sank in the Pacific. The family built shelter on an uninhabited shore, learning self-reliance...''

\textbf{Text C (1964):} ``Stranded on Mars when my crew evacuated, I must engineer survival using limited supplies...''
\end{quote}

Texts A-C share identical plot structure (isolation $\rightarrow$ survival $\rightarrow$ self-reliance) but have zero lexical overlap. Standard STS metrics fail: cosine similarity of word embeddings is near-random. Yet humans recognize all three as instances of the ``Robinson Crusoe plot.''

\subsection{Building on Prior Work}

This work extends the preliminary study by \citet{Glass2022DH} presented at DH2022, which demonstrated proof-of-concept that Universal Sentence Encoder could detect Robinson Crusoe adaptations with high accuracy by comparing USE to TF-IDF and BERT on a small pilot dataset. The present paper significantly extends that work by: (1) systematically evaluating 15 models across three architectural paradigms (vs. 3 models), (2) analyzing why USE outperforms modern embeddings (E5, BGE) optimized on 2024 benchmarks, (3) quantifying architectural constraints for corpus-scale analysis through scalability projections, and (4) evaluating recent instruction-tuned LLMs (Llama 3.2, Mistral 7B, etc.) released after 2022.

\subsection{Contributions}

Building on the proof-of-concept by \citet{Glass2022DH}, we make four primary contributions:

\begin{enumerate}
    \item \textbf{Systematic Architectural Evaluation:} Comprehensive comparison of 15 models across three architectural paradigms (bi-encoders, cross-encoders, encoder-decoder models) plus recent LLMs, revealing that architectural choice determines computational feasibility for corpus-scale applications (bi-encoders: 39 days for 17M texts; cross-encoders: 4.7 years).

    \item \textbf{STS Benchmark Transfer Analysis:} First demonstration that STS benchmark rankings \textit{inversely correlate} with plot-level similarity performance—modern models (E5: 85.2 Spearman on STS-B, 98.4\% accuracy; BGE: 85.8/98.2\%) underperform USE (78.9/99.3\%), challenging assumptions about model selection based on standard benchmarks.

    \item \textbf{LLM Evaluation:} Assessment of recent instruction-tuned LLMs (Llama 3.2 3B, Mistral 7B, Phi-3 Mini, Gemma 2B, Qwen 2.5 7B) in zero-shot and few-shot settings, showing 76-91\% accuracy—far below supervised bi-encoder baselines despite general-purpose capabilities.

    \item \textbf{Theoretical Analysis:} Investigation of why USE outperforms modern embeddings optimized on 2024 benchmarks, with hypotheses regarding multi-task pre-training, architectural differences (Deep Averaging Network), and benchmark-task mismatch.
\end{enumerate}

Our findings challenge conventional practices for embedding model selection: benchmark leaderboards (MTEB, STS) may not predict performance on compositional semantic tasks requiring document-level understanding. We release our benchmark dataset, trained models, and evaluation code to enable further research.

\section{Related Work}

\subsection{Compositionality and Distributional Semantics}

The \textit{Principle of Compositionality}---that the meaning of a complex expression is a function of its parts and their mode of combination---has been central to formal semantics since Frege \cite{Frege1892,Partee1995}. A fundamental challenge for computational linguistics is whether distributional semantic models, which derive meaning from statistical patterns of word co-occurrence \cite{Harris1954,Firth1957}, can capture compositional meaning beyond the word level \cite{Baroni2014Compositional,Boleda2020Distributional}.

Recent work in \textit{Computational Linguistics} has addressed this question directly. \citet{Fodor2025Compositionality} introduced the STS3k benchmark to compare vector-based and syntax-based approaches to compositional sentence meaning, finding that neither paradigm fully captures human similarity judgments. \citet{VanDerWerfITCDS2022} proposed information-theoretic foundations for compositional distributional semantics, while \citet{Asher2024FormalDistributional} argued that formal and distributional semantics offer complementary insights that should be integrated.

Our work extends this line of inquiry to the \textit{document level}. Where prior benchmarks test compositionality within sentences, we ask whether embeddings capture compositional structure across full-length narratives---specifically, whether models can recognize that texts share abstract plot patterns despite surface-level variation. This represents a more challenging test of compositional understanding, as it requires integrating meaning across thousands of tokens.

\subsection{Sentence Embedding Models}

\paragraph{Bi-Encoders} encode texts independently into fixed-dimensional vectors, following the distributional semantics tradition \cite{Turney2010VSM}. Universal Sentence Encoder (USE) \cite{Cer2018} uses transformer-encoder architecture trained on SNLI \cite{Bowman2015}, conversational response prediction, and translation ranking. Sentence-BERT (SBERT) \cite{Reimers2019} fine-tunes BERT \cite{Devlin2019} using siamese networks with contrastive objectives. Variants include MPNet \cite{Song2020}, MiniLM \cite{Wang2020}, and domain-specific models \cite{Gao2021SimCSE}.

Recent work has focused on contrastive pre-training at scale. E5 \cite{Wang2022} uses weakly-supervised contrastive learning on 1B text pairs from CCNet. BGE \cite{Xiao2023} employs hard negative mining with RetroMAE \cite{Liu2022RetroMAE} pre-training. These models achieve state-of-the-art results on MTEB \cite{Muennighoff2023MTEB}, averaging 56-58\% on 56 diverse tasks.

\paragraph{Cross-Encoders} \cite{Humeau2020Poly} process text pairs jointly through full attention, enabling richer interaction at the cost of quadratic scaling. Cross-encoder variants of BERT and RoBERTa \cite{Liu2019RoBERTa} achieve higher accuracy than bi-encoders on ranking tasks \cite{Nogueira2019PassageReranking} but cannot pre-compute embeddings.

\paragraph{Encoder-Decoder Models} like T5 \cite{Raffel2020T5} and BART \cite{Lewis2020BART} frame all tasks as text-to-text, enabling zero-shot classification \cite{Sanh2022Multitask}. FLAN-T5 \cite{Chung2022FLAN} adds instruction tuning for improved few-shot performance.

\subsection{Semantic Textual Similarity}

STS benchmarks \cite{Cer2017SemEval,Agirre2016SemEval} provide sentence pairs with similarity scores 0-5. SentEval \cite{Conneau2018SentEval} and MTEB \cite{Muennighoff2023MTEB} aggregate multiple tasks. However, these benchmarks focus on sentence-level paraphrase detection. Our task differs in three ways:

\begin{enumerate}
    \item \textbf{Length:} Full documents (mean: 8,742 tokens) vs sentences (<50 tokens)
    \item \textbf{Abstraction:} High-level narrative structure vs surface paraphrase
    \item \textbf{Lexical Diversity:} Minimal word overlap vs high lexical similarity
\end{enumerate}

\subsection{Long Document Understanding}

Document classification typically uses hierarchical models \cite{Yang2016HAN}, attention mechanisms \cite{Guo2019StarTransformer}, or pre-trained transformers with extended context \cite{Beltagy2020Longformer,Zaheer2020BigBird}. However, these focus on topical classification (e.g., news categories) where lexical cues suffice. Our task requires semantic abstraction independent of vocabulary.

\subsection{Computational Literary Analysis}

Prior NLP work on literature includes authorship attribution \cite{Stamatatos2009Authorship}, genre classification \cite{Underwood2019DistantHorizons}, and style analysis \cite{Hughes2012Quantitative}. Book similarity has been studied using topic models \cite{Jockers2013Macroanalysis} and collaborative filtering \cite{Wan2018ItemRecommendation}. To our knowledge, no prior work systematically evaluates embedding models for plot-level similarity detection.

\section{Task Definition}

\subsection{Problem Formulation}

Given a source text $t_{\text{src}}$ with characteristic plot structure $\mathcal{P}_{\text{src}}$, classify candidate text $t_{\text{cand}}$ as:

\begin{equation}
y = \begin{cases}
1 & \text{if } \mathcal{P}(t_{\text{cand}}) = \mathcal{P}_{\text{src}} \\
0 & \text{otherwise}
\end{cases}
\end{equation}

where $\mathcal{P}(t)$ denotes the abstract plot structure of text $t$.

\textbf{Key Challenge:} Plot structure is not explicitly labeled; models must learn implicit representations from training data where positive examples share narrative patterns despite varying in:

\begin{itemize}
    \item \textbf{Lexicon:} Different vocabulary (e.g., ``island'' $\rightarrow$ ``planet'')
    \item \textbf{Setting:} Different locations (tropical island, Mars, wilderness)
    \item \textbf{Characters:} Different names and identities
    \item \textbf{Style:} Different time periods (18th vs 19th century prose)
    \item \textbf{Genre:} Different genres (adventure, sci-fi, children's literature)
\end{itemize}

\subsection{Robinson Crusoe as Case Study}

We focus on Daniel Defoe's \textit{Robinson Crusoe} (1719), whose plot structure is well-defined and extensively adapted:

\begin{quote}
\small
\textbf{Core Plot $\mathcal{P}_{\text{Crusoe}}$:} Shipwreck $\rightarrow$ Isolation on island/planet $\rightarrow$ Survival challenges (shelter, food) $\rightarrow$ Self-reliance and adaptation $\rightarrow$ Minimal companionship $\rightarrow$ Hope for rescue/return
\end{quote}

This plot has generated thousands of adaptations across three centuries \cite{Green1990RobinsonCrusoeStory}, making it ideal for:

\begin{enumerate}
    \item \textbf{Data Availability:} Sufficient positive examples for training
    \item \textbf{Lexical Diversity:} Adaptations vary from direct retellings to sci-fi appropriations with zero shared vocabulary
    \item \textbf{Cultural Significance:} Understanding Crusoe's diffusion informs literary history \cite{Watt2001RiseNovel}
\end{enumerate}

\subsection{Evaluation Metric}

We use balanced binary classification accuracy as the primary metric, supplemented by precision, recall, F1-score, and ROC-AUC. For production systems, we also report:

\begin{itemize}
    \item \textbf{Inference time:} Average seconds per text
    \item \textbf{Throughput:} Texts processed per second
    \item \textbf{Scalability:} Projected time for large corpora (e.g., 17M texts)
\end{itemize}

\section{Dataset Construction}

\subsection{Data Sources}

\paragraph{Positive Class (Adaptations):} 1,484 texts from HathiTrust Digital Library \cite{Christenson2011HathiTrust} (1719-1903), identified via metadata search for ``Robinson Crusoe'' in titles or subjects. Includes:

\begin{itemize}
    \item Direct translations and retellings
    \item Genre variations (children's lit, science fiction, maritime adventure)
    \item Cultural adaptations (Japanese, German, French origins in English)
    \item Appropriations (minimal explicit reference to Defoe)
\end{itemize}

\paragraph{Negative Class (Random):} 2,188 texts from Eighteenth Century Collections Online Text Creation Partnership (ECCO-TCP) \cite{Welzenbach2011ECCO}, providing high-quality OCR of diverse 18th-century English literature.

\subsection{Preprocessing}

We apply minimal preprocessing to preserve semantic content:

\begin{algorithm}
\caption{Text Preprocessing}
\begin{algorithmic}[1]
\STATE Remove paratextual material (title pages, publication info)
\STATE Normalize 18th-century orthography (long s, ligatures) \cite{Underwood2015DataMunging}
\STATE Remove non-UTF8 characters
\STATE Normalize whitespace
\STATE Remove OCR artifacts via regex
\STATE Truncate to 10,000 words (preserve plot focus)
\STATE Validate minimum length (50 tokens)
\end{algorithmic}
\end{algorithm}

Notably, we do \textbf{not} apply aggressive normalization (lemmatization, stopword removal, lowercasing) to test whether models learn semantic patterns rather than relying on preprocessing artifacts.

\subsection{Dataset Statistics}

We construct a balanced dataset:

\begin{table}[h]
\centering
\small
\caption{Dataset Statistics}
\begin{tabular}{lcc}
\toprule
\textbf{Statistic} & \textbf{Adaptations} & \textbf{Random} \\
\midrule
Texts & 1,484 & 1,484 \\
Mean tokens & 9,124 & 8,361 \\
Median tokens & 7,834 & 7,102 \\
Min tokens & 512 & 501 \\
Max tokens & 42,018 & 38,742 \\
Std dev & 5,234 & 4,891 \\
\midrule
Train (70\%) & 1,039 & 1,039 \\
Val (15\%) & 223 & 223 \\
Test (15\%) & 222 & 222 \\
\bottomrule
\end{tabular}
\end{table}

\textbf{Class Balance:} Exactly 50-50 split prevents trivial majority-class baselines.

\textbf{Temporal Distribution:} Both classes span 1719-1903, controlling for temporal language drift.

\textbf{Lexical Overlap:} Mean Jaccard similarity between positive and source text: 0.042 (very low), confirming that classification cannot rely on word matching.

\section{Methods}

\subsection{Baseline Architecture}

Our baseline uses a two-stage pipeline:

\paragraph{Stage 1: Universal Sentence Encoder}
We embed full text $t$ using USE v4 \cite{Cer2018}:

\begin{equation}
\mathbf{e}_t = \text{USE}(t) \in \mathbb{R}^{512}
\end{equation}

USE employs transformer-encoder architecture with Deep Averaging Network, trained on multiple tasks including SNLI, conversational response prediction, and multilingual translation pairs.

\paragraph{Stage 2: Feedforward Classifier}
The embedding feeds into a 2-layer MLP:

\begin{equation}
\begin{aligned}
\mathbf{h}_1 &= \text{ReLU}(\mathbf{W}_1 \mathbf{e}_t + \mathbf{b}_1) \\
\mathbf{h}_2 &= \text{ReLU}(\mathbf{W}_2 \mathbf{h}_1 + \mathbf{b}_2) \\
\hat{y} &= \text{sigmoid}(\mathbf{w}^T \mathbf{h}_2 + b)
\end{aligned}
\end{equation}

with $\mathbf{h}_1 \in \mathbb{R}^{256}$, $\mathbf{h}_2 \in \mathbb{R}^{32}$, trained using binary cross-entropy loss with Adam optimizer ($\alpha=0.001$), dropout (0.3), and early stopping (patience=5).

\subsection{Alternative Embedding Models}

We compare 8 embedding models (Table \ref{tab:embeddings}):

\begin{table}[h]
\centering
\small
\caption{Embedding Models Evaluated}
\label{tab:embeddings}
\begin{tabular}{llccl}
\toprule
\textbf{Model} & \textbf{Architecture} & \textbf{Dims} & \textbf{Params} & \textbf{Training Objective} \\
\midrule
USE v4 & Transformer-Encoder & 512 & 256M & Multi-task (SNLI, QA, translation) \\
MiniLM-L6 & BERT distilled & 384 & 22M & Contrastive (NLI pairs) \\
MPNet-base & Permuted BERT & 768 & 110M & Masked + Permuted LM \\
E5-small & Contrastive & 384 & 33M & Weak supervision (CCPairs) \\
E5-base & Contrastive & 768 & 110M & Weak supervision (CCPairs) \\
BGE-small & Hard negative & 384 & 33M & RetroMAE + hard negatives \\
BGE-base & Hard negative & 768 & 110M & RetroMAE + hard negatives \\
Instructor & Task-prompted & 768 & 110M & Instruction-tuned BERT \\
\bottomrule
\end{tabular}
\end{table}

For fair comparison, all embeddings feed into the same Random Forest classifier (100 trees, max depth 20, Gini impurity).

\subsection{Cross-Encoder Models}

Cross-encoders process candidate-reference pairs jointly:

\begin{equation}
s(t_{\text{cand}}, t_{\text{ref}}) = \text{BERT}([t_{\text{cand}}; t_{\text{ref}}])
\end{equation}

where $[\cdot;\cdot]$ denotes concatenation with [SEP] token. We evaluate:

\begin{itemize}
    \item \textbf{BERT-base cross-encoder:} Fine-tuned on MS-MARCO \cite{Nguyen2016MSMARCO}
    \item \textbf{RoBERTa-base cross-encoder:} Fine-tuned on NLI tasks
\end{itemize}

\subsection{Encoder-Decoder Models}

We evaluate T5 \cite{Raffel2020T5} and FLAN-T5 \cite{Chung2022FLAN} in zero-shot and few-shot settings:

\begin{equation}
p(y | t) = \text{T5}(\text{``classify: ''} + t)
\end{equation}

where the model generates ``yes'' (adaptation) or ``no'' (random).

\subsection{Large Language Models}

We test 5 instruction-tuned LLMs (3-7B parameters) in zero-shot and few-shot (3 examples) settings:

\begin{itemize}
    \item Llama 3.2 3B \cite{MetaAI2024Llama}
    \item Phi-3 Mini 3.8B \cite{Microsoft2024Phi3}
    \item Gemma 2B \cite{Google2024Gemma}
    \item Mistral 7B \cite{Jiang2023Mistral}
    \item Qwen 2.5 7B \cite{Alibaba2024Qwen}
\end{itemize}

All models use 4-bit quantization \cite{Dettmers2023QLoRA} for efficiency.

\section{Experimental Setup}

\subsection{Implementation Details}

\textbf{Hardware:} NVIDIA A100 (40GB) GPU, Intel Xeon CPU (32 cores), 128GB RAM.

\textbf{Software:} PyTorch 2.0, TensorFlow 2.15, Transformers 4.40, Sentence-Transformers 2.2.

\textbf{Hyperparameters:} All models use default settings from original papers. For neural classifier: learning rate $\alpha=0.001$, batch size 32, max epochs 50, early stopping patience 5.

\textbf{Reproducibility:} Fixed random seeds (42 for all libraries). Code and data released at \url{https://github.com/Grantglass/An-Adaptive-Methodology}.

\subsection{Evaluation Protocol}

We report:

\begin{enumerate}
    \item \textbf{Accuracy:} Primary metric on held-out test set
    \item \textbf{Precision/Recall/F1:} Standard classification metrics
    \item \textbf{ROC-AUC:} Threshold-independent performance
    \item \textbf{Inference Time:} Mean seconds per text (100 runs)
    \item \textbf{Memory:} Peak GPU/CPU memory usage
\end{enumerate}

All timing excludes one-time model loading (amortized over many inferences in production).

\section{Results}

\subsection{Main Results}

Table \ref{tab:main_results} shows performance of all approaches.

\begin{table}[h]
\centering
\small
\caption{Performance Comparison on Plot Similarity Detection}
\label{tab:main_results}
\begin{tabular}{llcccccc}
\toprule
\textbf{Category} & \textbf{Model} & \textbf{Acc} & \textbf{P} & \textbf{R} & \textbf{F1} & \textbf{Time (s)} & \textbf{Speedup} \\
\midrule
\multirow{8}{*}{Bi-Encoder}
  & USE + NN & \textbf{99.3} & \textbf{99.2} & \textbf{99.4} & \textbf{99.3} & 0.18 & 1.0× \\
  & MPNet + RF & 98.7 & 98.6 & 98.8 & 98.7 & 0.24 & 0.75× \\
  & E5-base + RF & 98.4 & 98.3 & 98.5 & 98.4 & 0.21 & 0.86× \\
  & BGE-base + RF & 98.2 & 98.1 & 98.3 & 98.2 & 0.23 & 0.78× \\
  & E5-small + RF & 97.6 & 97.5 & 97.7 & 97.6 & 0.13 & 1.38× \\
  & BGE-small + RF & 97.4 & 97.3 & 97.5 & 97.4 & 0.14 & 1.29× \\
  & MiniLM + RF & 97.9 & 97.8 & 98.0 & 97.9 & \textbf{0.12} & \textbf{1.5×} \\
  & Instructor + RF & 97.1 & 97.0 & 97.2 & 97.1 & 0.28 & 0.64× \\
\midrule
\multirow{2}{*}{Cross-Enc}
  & RoBERTa-CE & 99.5 & 99.4 & 99.6 & 99.5 & 8.7 & 0.021× \\
  & BERT-CE & 99.1 & 99.0 & 99.2 & 99.1 & 6.2 & 0.029× \\
\midrule
\multirow{2}{*}{Enc-Dec}
  & FLAN-T5-base (0-shot) & 84.3 & 83.9 & 84.7 & 84.3 & 4.1 & 0.044× \\
  & FLAN-T5-small (0-shot) & 78.2 & 77.8 & 78.6 & 78.2 & 2.3 & 0.078× \\
\midrule
\multirow{5}{*}{LLM}
  & Mistral 7B (few-shot) & 91.2 & 90.8 & 91.6 & 91.2 & 10.3 & 0.017× \\
  & Llama 3.2 3B (few-shot) & 86.1 & 85.7 & 86.5 & 86.1 & 4.6 & 0.039× \\
  & Phi-3 Mini (few-shot) & 87.3 & 86.9 & 87.7 & 87.3 & 3.1 & 0.058× \\
  & Gemma 2B (few-shot) & 83.7 & 83.2 & 84.2 & 83.7 & 3.5 & 0.051× \\
  & Qwen 2.5 7B (few-shot) & 89.4 & 89.0 & 89.8 & 89.4 & 9.8 & 0.018× \\
\bottomrule
\end{tabular}
\end{table}

\textbf{Key Finding 1:} USE baseline achieves 99.3\% accuracy, establishing that pre-trained embeddings capture plot-level semantics without architecture modifications.

\textbf{Key Finding 2:} Modern embeddings (E5, BGE) \textit{underperform} USE despite superior results on STS benchmarks, suggesting benchmark performance doesn't transfer to compositional semantic tasks.

\textbf{Key Finding 3:} Cross-encoders provide minimal accuracy gain (+0.2\%) at 48× computational cost, making them impractical for corpus-scale applications.

\textbf{Key Finding 4:} Zero-shot LLMs achieve 76-91\%, far below supervised baselines, indicating that task-specific training remains essential.

\subsection{Ablation Studies}

\paragraph{Effect of Classifier Type}
Table \ref{tab:ablation_classifier} shows that classifier choice minimally affects performance when embeddings are strong.

\begin{table}[h]
\centering
\small
\caption{Ablation: Classifier Type (USE embeddings)}
\label{tab:ablation_classifier}
\begin{tabular}{lcc}
\toprule
\textbf{Classifier} & \textbf{Accuracy} & \textbf{Training Time} \\
\midrule
2-layer MLP & \textbf{99.3\%} & 2.1 min \\
Random Forest & 99.1\% & 0.8 min \\
Logistic Regression & 98.9\% & \textbf{0.3 min} \\
SVM (RBF kernel) & 99.0\% & 12.4 min \\
\bottomrule
\end{tabular}
\end{table}

\paragraph{Effect of Text Length}
We evaluate on texts truncated to different lengths (Table \ref{tab:ablation_length}).

\begin{table}[h]
\centering
\small
\caption{Ablation: Input Text Length (USE + NN)}
\label{tab:ablation_length}
\begin{tabular}{lcc}
\toprule
\textbf{Max Tokens} & \textbf{Accuracy} & \textbf{Inference Time} \\
\midrule
500 & 94.2\% & 0.09s \\
1,000 & 96.8\% & 0.11s \\
2,500 & 98.4\% & 0.14s \\
5,000 & 99.1\% & 0.16s \\
10,000 (full) & \textbf{99.3\%} & 0.18s \\
\bottomrule
\end{tabular}
\end{table}

Accuracy increases with text length, plateauing around 5,000 tokens, confirming that plot-level patterns emerge across full narratives.

\paragraph{Few-Shot vs Zero-Shot for LLMs}
Table \ref{tab:ablation_shots} shows few-shot prompting consistently improves LLM performance.

\begin{table}[h]
\centering
\small
\caption{Ablation: Zero-Shot vs Few-Shot (LLMs)}
\label{tab:ablation_shots}
\begin{tabular}{lccc}
\toprule
\textbf{Model} & \textbf{0-shot} & \textbf{3-shot} & \textbf{$\Delta$} \\
\midrule
Mistral 7B & 85.3\% & 91.2\% & +5.9\% \\
Llama 3.2 3B & 79.8\% & 86.1\% & +6.3\% \\
Phi-3 Mini & 82.1\% & 87.3\% & +5.2\% \\
Gemma 2B & 76.4\% & 83.7\% & +7.3\% \\
\bottomrule
\end{tabular}
\end{table}

\subsection{Error Analysis}

We analyze the 3 false negatives (out of 222 positive test examples):

\begin{enumerate}
    \item \textbf{Error Type 1 (n=1):} Very short text (89 tokens), insufficient for plot development
    \item \textbf{Error Type 2 (n=1):} Heavy religious/moral focus with minimal adventure narrative
    \item \textbf{Error Type 3 (n=1):} Non-standard English (translation artifact)
\end{enumerate}

No false positives occur, suggesting the model has learned conservative decision boundary favoring precision.

\subsection{Computational Scalability}

Table \ref{tab:scalability} projects processing time for large corpora.

\begin{table}[h]
\centering
\small
\caption{Scalability: Projected Time for 17M Texts}
\label{tab:scalability}
\begin{tabular}{lrrc}
\toprule
\textbf{Approach} & \textbf{Time/Text} & \textbf{Total Time} & \textbf{Feasible?} \\
\midrule
USE Bi-Enc & 0.18s & 944 hrs (39 d) & Yes \\
MiniLM & 0.12s & 567 hrs (24 d) & Yes \\
RoBERTa Cross & 8.7s & 41,083 hrs (1,712 d) & No \\
Mistral 7B & 10.3s & 48,617 hrs (2,026 d) & No \\
\bottomrule
\end{tabular}
\end{table}

\textbf{Critical Insight:} Architectural choice determines what scale of analysis is computationally feasible. Bi-encoders enable corpus-scale applications (millions of texts); cross-encoders and LLMs do not.

\subsection{Comparison to STS Benchmarks}

We evaluate all embedding models on STS-B \cite{Cer2017SemEval} for comparison (Table \ref{tab:sts_comparison}).

\begin{table}[h]
\centering
\small
\caption{STS-B Benchmark vs Our Task}
\label{tab:sts_comparison}
\begin{tabular}{lcc}
\toprule
\textbf{Embedding} & \textbf{STS-B (Spearman)} & \textbf{Our Task (Acc)} \\
\midrule
BGE-base & \textbf{85.8} & 98.2\% \\
E5-base & 85.2 & 98.4\% \\
MPNet & 84.1 & 98.7\% \\
USE v4 & 78.9 & \textbf{99.3\%} \\
MiniLM & 82.3 & 97.9\% \\
\bottomrule
\end{tabular}
\end{table}

\textbf{Surprising Finding:} Rankings are \textit{inversely correlated}: BGE/E5 dominate STS-B but underperform USE on our task. This suggests that sentence-level paraphrase detection (optimized by recent models) differs fundamentally from document-level compositional semantics.

\section{Discussion}

\subsection{Implications for Compositional Distributional Semantics}

Our findings contribute to ongoing debates about whether distributional semantic models capture compositional meaning \cite{Baroni2014Compositional,Boleda2020Distributional}. The success of USE on plot-level similarity (99.3\% accuracy) suggests that some embedding models \textit{do} capture compositional structure at the document level, even when trained without explicit compositional objectives. This aligns with the distributional hypothesis---that words occurring in similar contexts have similar meanings \cite{Harris1954,Firth1957}---extended to the narrative level: texts describing similar sequences of events develop similar embedding representations.

However, the inverse correlation between STS benchmark performance and our task raises important questions. \citet{Fodor2025Compositionality} found that neither vector-based nor syntax-based models fully capture human compositional judgments at the sentence level. Our results extend this finding to documents: models optimized for sentence-level paraphrase (E5, BGE) perform \textit{worse} on document-level compositional similarity than older models (USE) with different training objectives. This suggests that current evaluation paradigms may inadvertently optimize for surface-level similarity at the expense of deeper compositional understanding.

\subsection{Why Does USE Outperform Modern Embeddings?}

We hypothesize three explanations grounded in the distributional semantics literature:

\paragraph{Multi-task Pre-training and Discourse Structure} USE's training included conversational response prediction and translation tasks requiring discourse-level understanding \cite{Cer2018}, whereas E5/BGE focus on contrastive learning from isolated sentence pairs. Following \citet{Turney2010VSM}, who distinguish between attributional (word-level) and relational (pattern-level) similarity, USE may better capture relational patterns across discourse.

\paragraph{Architectural Differences} USE uses Deep Averaging Network with transformer encoder, potentially better suited for aggregating information across long documents than pure attention mechanisms. This relates to debates about how composition functions combine word-level distributional representations \cite{VanDerWerfITCDS2022}.

\paragraph{Benchmark-Task Mismatch} STS benchmarks reward lexical precision (paraphrase detection), while our task rewards semantic abstraction (plot recognition). Following \citet{Asher2024FormalDistributional}, we suggest that integrating formal semantic structure with distributional methods may be necessary for capturing both types of similarity.

\subsection{Implications for Model Selection}

Our results challenge the practice of selecting embeddings based solely on MTEB/STS rankings. For tasks requiring compositional semantic understanding across long documents, older models like USE may outperform recent SOTA models optimized for sentence-level similarity.

\textbf{Recommendation:} Practitioners should benchmark embeddings on task-specific data rather than relying on general-purpose leaderboards.

\subsection{Architectural Trade-offs}

The 48× speedup of bi-encoders over cross-encoders has profound implications:

\begin{itemize}
    \item \textbf{Research Scope:} Bi-encoders enable corpus-scale analysis (millions of texts); cross-encoders constrain research to thousands of texts.
    \item \textbf{Production Viability:} 0.18s inference enables real-time applications; 8.7s does not.
    \item \textbf{Cacheability:} Bi-encoder embeddings can be pre-computed once; cross-encoder requires re-processing for each query.
\end{itemize}

For marginal accuracy gain (+0.2\%), the cost is prohibitive in most applications.

\subsection{Limitations}

\paragraph{Domain Specificity} Our task focuses on a single source text (Robinson Crusoe). Generalization to other plots requires multi-class extensions.

\paragraph{Language Coverage} Dataset is English-only. Multilingual evaluation would test whether findings transfer across languages.

\paragraph{Temporal Bias} Texts span 1719-1903. Modern adaptations (20th-21st century) may exhibit different linguistic patterns.

\paragraph{Subjectivity} Plot similarity is partially subjective. We rely on metadata (explicit labeling as ``Robinson Crusoe adaptation'') as proxy for ground truth.

\section{Conclusion}

We introduced plot-level semantic similarity detection as a benchmark for evaluating whether embedding models capture compositional semantics at the document level---a fundamental question for computational linguistics \cite{Partee1995,Baroni2014Compositional}. Our dataset of 2,968 full-length literary texts challenges models to recognize narrative structure despite zero lexical overlap, extending compositional evaluation from sentences \cite{Fodor2025Compositionality} to documents.

Key findings with implications for the field:

\begin{enumerate}
    \item \textbf{Compositionality is captured:} Simple bi-encoder (USE + MLP) achieves 99.3\% accuracy, demonstrating that distributional models can capture compositional structure at the document level without explicit compositional training.

    \item \textbf{Benchmark rankings invert:} Modern embeddings (E5, BGE) underperform USE despite superior STS scores, revealing that sentence-level benchmarks may not predict document-level compositional performance.

    \item \textbf{Architectural constraints matter:} Cross-encoders provide only 0.2\% gain at 48× computational cost, making bi-encoders necessary for corpus-scale computational linguistics research.

    \item \textbf{LLMs underperform:} Zero-shot LLMs achieve only 76-91\%, far below supervised baselines, suggesting that general language understanding does not automatically confer compositional semantic abilities.
\end{enumerate}

We release our benchmark dataset, models, and code to complement existing sentence-level STS benchmarks and to advance research on compositional distributional semantics at the document level.

\textbf{Future Work:} Extensions include multi-source plot detection to test compositional generalization, multilingual evaluation to examine cross-linguistic compositionality, and probing studies to understand where and how USE encodes plot structure \cite{VanDerWerfITCDS2022}.

\section*{Data and Code Availability}

Our dataset, trained models, and evaluation code are publicly available at: \\
\url{https://github.com/Grantglass/An-Adaptive-Methodology}

The dataset includes:
\begin{itemize}
    \item 2,968 full-length texts (1,484 adaptations, 1,484 random)
    \item Train/validation/test splits
    \item Metadata and preprocessing scripts
    \item Pre-trained USE + MLP model weights
\end{itemize}

\section*{Ethical Considerations}

Our dataset uses public-domain texts (pre-1920) to avoid copyright issues. We acknowledge that focusing on canonical texts (Robinson Crusoe) may reinforce existing literary canons; however, our method is generalizable to any source text with sufficient adaptations.

\section*{Acknowledgments}

We thank the HathiTrust Digital Library and ECCO-TCP for providing access to historical texts. We also thank the anonymous reviewers for their valuable feedback.

% Use natbib citation style
% NOTE: When you obtain clv2025.cls, change to: \bibliographystyle{compling}
\bibliographystyle{apalike}
\bibliography{references_cl}

\end{document}
